% ─────────────────────────────────────────────────────────────────────────────
%  Sección 4.X  Diseño Preliminar del Módulo de Análisis OCR y NLP
%  NOTA PARA EL EQUIPO: Esta sección corresponde al diseño de la iteración
%  siguiente (TT2). Se presenta a nivel arquitectónico y conceptual, dejando
%  abiertas las decisiones de implementación que serán tomadas conforme al
%  proceso de desarrollo en espiral adoptado por el proyecto.
% ─────────────────────────────────────────────────────────────────────────────

\subsection{Diseño Preliminar del Módulo de Análisis OCR y NLP}

\textit{El presente apartado constituye un diseño preliminar enmarcado en
la metodología de desarrollo en espiral adoptada por el proyecto. Dado que
la implementación completa de este módulo corresponde a la siguiente
iteración (TT2), las decisiones que se presentan a continuación tienen
carácter arquitectónico y conceptual: establecen las interfaces, los
contratos de datos y la posición del módulo dentro del sistema, pero dejan
abiertas las decisiones de modelo, hiperparámetros y bibliotecas específicas,
las cuales serán evaluadas y refinadas con base en las pruebas de la iteración
actual. Esto es consistente con el enfoque en espiral, cuya característica
distintiva es precisamente diferir las decisiones de alto riesgo técnico hasta
que se cuente con evidencia empírica suficiente para tomarlas.}

\vspace{0.5em}

El módulo de análisis OCR y NLP constituye el núcleo de procesamiento
inteligente de la plataforma. Su función es transformar la imagen de escritura
manuscrita capturada por el alumno en dos conjuntos de resultados cuantitativos
y cualitativos: las métricas caligráficas (alineación, tamaño, espaciado e
inclinación) y el análisis ortográfico del texto detectado. Los fundamentos
teóricos de ambas tecnologías han sido descritos en las Secciones~2.3 y~2.4
del presente documento; esta sección se limita a su diseño de integración
dentro de la arquitectura del sistema.

La arquitectura de este módulo sigue el patrón MVC establecido para el
sistema y se divide en tres componentes principales:

\begin{enumerate}
    \item \textbf{Capa de Procesamiento (Back-End):} Un controlador dedicado
    \texttt{analisis.py} recibirá como entrada la imagen en formato Base64
    almacenada en la tabla \texttt{Intentos} y orquestará la ejecución
    secuencial del sub-módulo OCR y del sub-módulo NLP. Este controlador
    será invocado de forma \textbf{diferida} respecto al registro del intento,
    es decir, el alumno recibe confirmación inmediata de envío y el
    procesamiento ocurre de manera asíncrona en segundo plano, evitando
    tiempos de espera perceptibles en la interfaz.

    \item \textbf{Sub-módulo OCR:} Responsable de preprocesar la imagen y
    extraer el texto manuscrito en formato digital, así como de calcular las
    métricas caligráficas del trazo.

    \item \textbf{Sub-módulo NLP:} Responsable de analizar el texto producido
    por el OCR, detectar errores ortográficos y generar sugerencias de
    corrección contextualizadas al nivel léxico de tercer y cuarto grado de
    primaria.
\end{enumerate}

% ─────────────────────────────────────────────────────────
\subsubsection{Diseño del Flujo de Procesamiento}

El flujo general del módulo, desde la captura del intento hasta el
almacenamiento de resultados, se describe a continuación:

\begin{enumerate}
    \item El alumno envía su intento mediante el endpoint
    \texttt{/api/intentos/registrar}. El Back-End almacena la imagen en
    Base64 en la tabla \texttt{Intentos} con los campos
    \texttt{texto\_detectado\_ocr} y los registros de análisis en estado
    nulo, confirmando el envío al cliente de forma inmediata.

    \item Un proceso en segundo plano —cuyo mecanismo de programación
    será definido en TT2, ya sea mediante un scheduler como el ya
    utilizado en \texttt{cerrar\_ejercicio.py}, una cola de tareas o una
    llamada asíncrona— detecta el nuevo intento pendiente de análisis.

    \item El \textbf{sub-módulo OCR} recibe la imagen en Base64, la
    decodifica y ejecuta la siguiente secuencia de procesamiento:
    \begin{enumerate}
        \item \textbf{Preprocesamiento:} conversión a escala de grises,
        binarización (método de Otsu \cite{otsu1979threshold}),
        eliminación de ruido y corrección de inclinación
        (\textit{deskewing}), conforme a las técnicas descritas en la
        Sección~2.3.
        \item \textbf{Extracción de métricas caligráficas:} cálculo de
        los cuatro parámetros definidos en el sistema:
        \texttt{alineacion\_score}, \texttt{tamano\_letra\_score},
        \texttt{espaciado\_score} e \texttt{inclinacion\_score}. La
        metodología exacta de cálculo será determinada en TT2.
        \item \textbf{Reconocimiento de texto:} aplicación del modelo de
        reconocimiento para obtener el texto digital. La elección del
        modelo (TrOCR, Tesseract u otro) queda abierta a la iteración
        siguiente, sujeta a las pruebas de precisión sobre escritura
        infantil descritas en el objetivo específico OE-03.
    \end{enumerate}

    \item El texto reconocido se escribe en el campo
    \texttt{texto\_detectado\_ocr} de la tabla \texttt{Intentos} y los
    scores caligráficos se persisten en la tabla
    \texttt{Analisis\_Caligrafico}.

    \item El \textbf{sub-módulo NLP} recibe el texto digitalizado y
    ejecuta:
    \begin{enumerate}
        \item \textbf{Tokenización:} división del texto en unidades
        léxicas individuales.
        \item \textbf{Detección de errores ortográficos:} comparación de
        cada token contra el modelo lingüístico, considerando el nivel
        cognitivo y vocabulario de tercer y cuarto grado de primaria.
        \item \textbf{Generación de sugerencias:} producción de
        correcciones contextualizadas mediante algoritmos de similitud
        o modelos estadísticos, conforme a lo descrito en la
        Sección~2.4.
    \end{enumerate}

    \item Los resultados del análisis ortográfico se persisten en la tabla
    \texttt{Analisis\_Ortografico}: el \texttt{ortografia\_score}, la
    \texttt{cantidad\_errores} y el campo \texttt{sugerencias\_json} con
    la estructura definida en la Sección~4.2.2.

    \item Con ambos análisis completados, el proceso actualiza el registro
    de \texttt{Progreso\_Alumno} mediante el trigger de base de datos
    descrito en la Sección~4.4, reflejando el nuevo promedio acumulado
    del alumno.
\end{enumerate}

% ─────────────────────────────────────────────────────────
\subsubsection{Contratos de Datos entre Módulos}

Dado que el diseño de integración se establece en TT1 para orientar el
desarrollo de TT2, se definen a continuación los contratos de datos que
las interfaces entre módulos deben respetar. Estos contratos son fijos
y no están sujetos a cambio en la iteración siguiente, pues representan
el acuerdo entre la base de datos ya implementada y el módulo por
desarrollar.

\paragraph{Entrada del sub-módulo OCR}

El sub-módulo OCR recibirá la imagen codificada en Base64 directamente
del campo \texttt{imagen\_codificada} de la tabla \texttt{Intentos},
identificada por su \texttt{id\_intento}. No se requiere ningún endpoint
adicional, ya que el proceso opera sobre datos ya persistidos.

\textbf{\textit{Salida del sub-módulo OCR hacia la base de datos}}

Actualización del campo \texttt{texto\_detectado\_ocr} en la tabla \texttt{Intentos}:

\begin{lstlisting}[language=json, caption={Estructura de resultados del sub-módulo OCR hacia la BD}, label={lst:ocr-output}]
{
  "texto_detectado_ocr": "string (texto reconocido)"
}

Registro a insertar en la tabla \texttt{Analisis\_Caligrafico}:

\begin{lstlisting}[language=json]
{
  "id_intento":          <integer>,
  "alineacion_score":    <decimal(5,2)>,
  "tamano_letra_score":  <decimal(5,2)>,
  "espaciado_score":     <decimal(5,2)>,
  "inclinacion_score":   <decimal(5,2)>,
  "observaciones":       "string | null"
}
\end{lstlisting}

\textbf{\textit{Salida del sub-módulo NLP hacia la base de datos}}

Registro a insertar en la tabla \texttt{Analisis\_Ortografico}:

\begin{lstlisting}[language=json, caption={Estructura de resultados del sub-módulo NLP hacia la BD}, label={lst:nlp-output}]
{
  "id_intento":        <integer>,
  "cantidad_errores":  <integer>,
  "ortografia_score":  <decimal(5,2)>,
  "errores_detectados": "string | null",
  "sugerencias_json": [
    {
      "palabra_incorrecta": "string",
      "posicion":           <integer>,
      "sugerencia":         "string",
      "contexto":           "string"
    }
  ]
}
\end{lstlisting}

El campo \texttt{sugerencias\_json} es validado al momento de la
inserción mediante la restricción \texttt{CK\_sugerencias\_json}
definida en el esquema de base de datos, la cual verifica que el
contenido sea JSON válido mediante la función \texttt{ISJSON()} de
SQL Server.

% ─────────────────────────────────────────────────────────
\subsubsection{Diseño del Endpoint de Consulta de Resultados}

Para que el Front-End pueda presentar los resultados del análisis al
alumno y al tutor, se contempla el siguiente endpoint en el controlador
\texttt{analisis.py}:

\paragraph{Endpoint: Consulta de Análisis por Intento}

El endpoint \texttt{/api/analisis/\{id\_intento\}} operará mediante el
método \texttt{GET} y requerirá autenticación JWT con rol de tutor o
alumno. Su propósito es retornar de forma consolidada los resultados
del análisis OCR y NLP asociados a un intento específico.

\begin{lstlisting}[language=json, caption={Estructura de respuesta del endpoint de análisis}, label={lst:analisis-response}]
{
  "id_intento":          <integer>,
  "texto_detectado_ocr": "string | null",
  "caligrafia": {
    "alineacion_score":   <decimal>,
    "tamano_letra_score": <decimal>,
    "espaciado_score":    <decimal>,
    "inclinacion_score":  <decimal>,
    "observaciones":      "string | null"
  },
  "ortografia": {
    "ortografia_score":   <decimal>,
    "cantidad_errores":   <integer>,
    "sugerencias": [
      {
        "palabra_incorrecta": "string",
        "posicion":           <integer>,
        "sugerencia":         "string",
        "contexto":           "string"
      }
    ]
  }
}
\end{lstlisting}

Si el análisis aún no ha sido procesado (campos nulos en la BD), el
endpoint retornará los campos de \texttt{caligrafia} y
\texttt{ortografia} como \texttt{null}, indicando al cliente que el
procesamiento se encuentra pendiente.

% ─────────────────────────────────────────────────────────
\subsubsection{Decisiones Abiertas para TT2}

En concordancia con el desarrollo en espiral, las siguientes decisiones
técnicas de alto riesgo se posponen deliberadamente a la siguiente
iteración, donde serán evaluadas mediante pruebas controladas:

\begin{itemize}
    \item \textbf{Modelo OCR:} la elección entre TrOCR, Tesseract
    \cite{smith2007overview} u otro motor queda condicionada a las
    pruebas de precisión sobre muestras de escritura infantil en letra
    de molde (tipografía Century) \cite{troncoso2009sindrome}, conforme al objetivo específico OE-03
    que establece una precisión mínima del 85\%.

    \item \textbf{Metodología de métricas caligráficas:} la definición
    de los algoritmos de cálculo de \texttt{alineacion\_score},
    \texttt{tamano\_letra\_score}, \texttt{espaciado\_score} e
    \texttt{inclinacion\_score} depende del análisis de las
    características del trazo reconocido por el motor OCR elegido.

    \item \textbf{Modelo NLP:} la selección entre un corrector basado en
    diccionario, distancia de Levenshtein o un modelo estadístico más
    avanzado \cite{jurafsky2023speech} queda sujeta a la evaluación de
    su desempeño con el vocabulario del nivel educativo objetivo.

    \item \textbf{Mecanismo de ejecución diferida:} la decisión de
    implementar el procesamiento asíncrono mediante un scheduler,
    una cola de tareas (como Celery) o corutinas de FastAPI queda
    abierta a la iteración siguiente, donde se evaluará el impacto
    en el rendimiento del servidor.
\end{itemize}

Lo que sí queda definido y es inmutable para TT2 son los contratos de
datos descritos en la Sección~4.X.2 y el esquema de base de datos
ya implementado, que garantizan la integración sin fricciones del
módulo una vez desarrollado.